% =====================================================================
% Template ANPEd — Trabalho Completo / Pôster
% Para submissões à Reunião Nacional da ANPEd e GTs
%
% NOTA: Cada GT da ANPEd pode ter ajustes específicos.
%       Sempre confira as normas do GT em https://www.anped.org.br
%
% Compilação: pdflatex doc.tex (3x para sumário/referências)
% =====================================================================

\documentclass[12pt,a4paper]{article}

% --- Pacotes essenciais
\usepackage[utf8]{inputenc}
\usepackage[T1]{fontenc}
\usepackage[brazil]{babel}
\usepackage{lmodern}
\usepackage{microtype}

% --- Margens (padrão ANPEd: 3-2-3-2 cm, Times/Arial 12, espaçamento 1,5)
\usepackage[a4paper,top=3cm,bottom=2cm,left=3cm,right=2cm]{geometry}

% --- Espaçamento 1,5
\usepackage{setspace}
\onehalfspacing

% --- Citações ABNT (autor-data)
\usepackage[alf,abnt-emphasize=bf,abnt-etal-list=0,abnt-etal-cite=3]{abntex2cite}

% --- Hiperlinks
\usepackage[colorlinks=true,citecolor=black,linkcolor=black,urlcolor=blue]{hyperref}

% --- Tabelas e figuras
\usepackage{graphicx}
\usepackage{booktabs}
\usepackage{caption}
\captionsetup[table]{position=above}
\captionsetup[figure]{position=below}

% --- Recuo de parágrafo (1,25 cm)
\setlength{\parindent}{1.25cm}
\setlength{\parskip}{0pt}

% --- Comandos auxiliares
\newcommand{\titulodotrabalho}{TÍTULO DO TRABALHO\\\large Subtítulo, se houver}
\newcommand{\gtanped}{GT XX -- Nome do Grupo de Trabalho}
\newcommand{\nomedoautor}{Nome Sobrenome -- Sigla IES}

% =====================================================================
\begin{document}

% --- Cabeçalho da ANPEd (sem identificação na primeira versão; só após aprovação)
\begin{center}
{\large \textbf{\titulodotrabalho}}\\[0.5cm]
\textbf{\gtanped}\\[0.3cm]
\nomedoautor\\
\href{mailto:email@uea.edu.br}{email@uea.edu.br}\\
% --- Apoio financeiro, se houver:
% Bolsa CAPES / CNPq nº XXXXX / Edital FAPEAM nº XX
\end{center}

\vspace{0.5cm}

% =====================================================================
% RESUMO + PALAVRAS-CHAVE (NBR 6028)
% =====================================================================
\noindent\textbf{Resumo:} A presente pesquisa investigou [tema] no
contexto de [recorte], adotando [tipo de pesquisa] como abordagem
metodológica. Foram coletados [dados] mediante [instrumento de coleta],
com posterior análise por meio de [método de análise]. Os resultados
indicaram que [achado 1], além de [achado 2]. Conclui-se que
[contribuição original], o que sinaliza para [implicações]. As
limitações deste estudo incluem [limitação principal], sugerindo-se,
para pesquisas futuras, [agenda].

\vspace{0.3cm}

\noindent\textbf{Palavras-chave:} termo 1; termo 2; termo 3; termo 4;
termo 5.

\vspace{0.5cm}

% =====================================================================
% INTRODUÇÃO
% =====================================================================
\section*{Introdução}

[Parágrafo 1 -- Contextualização do tema]

[Parágrafo 2 -- Lacuna na literatura, com 3-5 referências brasileiras]

[Parágrafo 3 -- Pergunta de pesquisa + objetivos: geral + 2-4 específicos]

[Parágrafo 4 -- Justificativa em 3 dimensões: teórica, prática, social]

[Parágrafo 5 -- Estrutura do trabalho: ``Após esta introdução, a seção
2 apresenta...'']

% =====================================================================
% REFERENCIAL TEÓRICO
% =====================================================================
\section*{Referencial Teórico}

\subsection*{[Conceito-chave 1]}

[Definição clássica + autor canônico]. \cite[p.~XX]{autor1}

[Evolução / debates contemporâneos]. \cite{autor2,autor3}

[Aplicação ao tema da pesquisa].

\subsection*{[Conceito-chave 2]}

[Estrutura semelhante]

% =====================================================================
% METODOLOGIA
% =====================================================================
\section*{Metodologia}

A pesquisa caracteriza-se, quanto à natureza, como
[básica/aplicada]. Quanto à abordagem, classifica-se como
[qualitativa/quantitativa/mista]. Em relação aos objetivos,
trata-se de pesquisa [exploratória/descritiva/explicativa]. Os
procedimentos adotados configuram [estudo de caso/levantamento/
análise documental/etnografia/pesquisa-ação]
\citep{gil2017,marconi2017}.

% --- Sujeitos / corpus
[Quem participou? Quantos? Critérios de seleção?]

% --- Coleta
[Que instrumento? Quando? Onde?]

% --- Análise
A análise foi realizada conforme proposta de \citeonline{bardin2011}, em
três fases: pré-análise, exploração do material e tratamento dos
resultados.

% --- Aspectos éticos
A pesquisa foi aprovada pelo Comitê de Ética em Pesquisa da [IES],
parecer nº XXXXX, CAAE XXXXX.

% =====================================================================
% RESULTADOS E DISCUSSÃO
% =====================================================================
\section*{Resultados e Discussão}

\subsection*{[Eixo 1, alinhado com objetivo específico 1]}

Constatou-se que [achado], conforme demonstrado na Tabela~\ref{tab:perfil}.

\begin{table}[htbp]
\centering
\caption{Perfil dos participantes (n=15)}
\label{tab:perfil}
\begin{tabular}{lcc}
\toprule
\textbf{Variável} & \textbf{n} & \textbf{\%} \\
\midrule
Gênero -- Feminino & 9 & 60{,}0 \\
Gênero -- Masculino & 6 & 40{,}0 \\
\bottomrule
\end{tabular}
\\\smallskip
\footnotesize Fonte: dados da pesquisa (2026).
\end{table}

[Interpretação do achado + diálogo com literatura]

\subsection*{[Eixo 2]}

[...]

% =====================================================================
% CONSIDERAÇÕES FINAIS
% =====================================================================
\section*{Considerações Finais}

[Recapitulação da pergunta]. [Síntese dos achados]. [Contribuição
em 2-3 dimensões]. [Limitações]. [Agenda futura].

% =====================================================================
% REFERÊNCIAS
% =====================================================================
\renewcommand{\refname}{Referências}
\bibliography{referencias}

\end{document}
